\documentclass[11pt]{article}
\usepackage{fontspec}
\setmainfont{Times New Roman}
\setsansfont{Arial}
\setmonofont{Consolas}
\usepackage{amsmath,amssymb,mathtools}
\usepackage{geometry}
\usepackage{microtype}
\newcommand{\mo}{\mathfrak{o}}
\newcommand{\mO}{\mathfrak{O}}
\newcommand{\mT}{\mathfrak{T}}
\newcommand{\mR}{\mathfrak{R}}
\newcommand{\mq}{\mathfrak{q}}
\newcommand{\ma}{\mathfrak{a}}
\newcommand{\mb}{\mathfrak{b}}
\newcommand{\mc}{\mathfrak{c}}
\newcommand{\mpideal}{\mathfrak{p}}
\geometry{a4paper,margin=2.35cm}
\setlength{\parindent}{1.25em}
\setlength{\parskip}{0pt}
\makeatletter
\renewcommand\footnotesize{\@setfontsize\footnotesize{7.7}{8.7}\selectfont}
\makeatother
\emergencystretch=100em
\allowdisplaybreaks
\sloppy
\hbadness=10000
\vbadness=10000
\hfuzz=2pt
\vfuzz=10pt
\raggedbottom

\begin{document}

% Unit: Paper31 Section08 Entry06 multiplication-ring quotient localization v001. Direct Interslavic Latin translation from clean RA34 line 439 / segment P31-S0099, with scan-visible final equality/proof sentence restored from printed page 103 / PDF page 22.
\setcounter{footnote}{37}

\noindent Te same razmysljenja važe takože za multiplikacijsko kolco $\mo$. Osoblivo, ako $\mpideal$ jest prosty ideal, togda $\mo_\mpideal$ imaje samo jedin prosty ideal, različny od nulovogo i jediničnogo ideala; zaradi izomorfizma kolc klasov ostatkov kolco $\mo_\mpideal$, kako i $\mo$, jest multiplikacijsko kolco, i po 1 togda $\mo_\mpideal$ jest kolco glavnyh idealov. Element, ktory poraždaje prosty ideal izvedeny iz $\mpideal$, stava prostym elementom $p$ iz $\mo_\mpideal$. Razširjeny ideal libovoljnogo ideala $\mc$ iz $\mo$ jest porodženy primarnoju komponentoju $\mc$, prinadležečoju k $\mpideal$; zato on jest ravny $(p)^\alpha$ abo jediničnomu idealu, soglasno tomu, vhodi li $\mpideal$ v $\mc$ -- i to v $\alpha$-toj stepeni -- abo ne vhodi v $\mc$. Iz togo još slěduje: ako razširjenja dvuh idealov $\mb$ i $\mc$ iz $\mo$ sovpadajut v vsih kolcah kvocientov $\mo_\mpideal$, togda $\mb$ i $\mc$ sut identične, bo oni sut dělimy tymi samymi prostymi idealami $\mpideal$ i v tyh samyh stepenjah.

\end{document}
